\documentclass{article}
\usepackage[utf8]{inputenc}
\usepackage[spanish]{babel}
\usepackage{amsmath}
\usepackage{amssymb}
\usepackage{amsthm}
\usepackage{geometry}

\geometry{margin=2.5cm}

\theoremstyle{definition}
\newtheorem{definition}{Definición}
\newtheorem{lemma}{Lema}

\title{Cálculo de la Desviación Estándar de la Distribución Exponencial}
\author{Juan Felipe Rodríguez G. (20261595004)}
\date{\today}

\begin{document}
    \maketitle

    \section*{Contexto del Problema}

    Sea $X$ una variable aleatoria continua que sigue una distribución exponencial con parámetro de tasa $\lambda > 0$. El objetivo de este documento es calcular su desviación estándar $\sigma_X$, desarrollando de forma completa cada momento necesario y, al final, interpretar el significado del resultado obtenido.

    \section*{Fundamentos Teóricos}

    \begin{definition}[Distribución exponencial]
    Decimos que $X \sim \mathrm{Exp}(\lambda)$, con $\lambda > 0$, si su función de densidad de probabilidad es:
    \begin{equation} \label{eq:pdf}
        f(x) =
        \begin{cases}
            \lambda e^{-\lambda x}, & x \geq 0, \\[4pt]
            0,                       & x < 0.
        \end{cases}
    \end{equation}
    \end{definition}

    \begin{definition}[Varianza y desviación estándar]
    La varianza de una variable aleatoria $X$ se define mediante la identidad de König-Huygens:
    \begin{equation} \label{eq:varianza}
        \operatorname{Var}(X) = E[X^2] - \left(E[X]\right)^2,
    \end{equation}
    donde $E[X]$ es el primer momento (esperanza) y $E[X^2]$ es el segundo momento. La desviación estándar es la raíz cuadrada no negativa de la varianza:
    \begin{equation} \label{eq:sigma}
        \sigma_X = \sqrt{\operatorname{Var}(X)}.
    \end{equation}
    \end{definition}

    \begin{definition}[Integración por partes]
    Para funciones diferenciables $u(x)$ y $v(x)$, se cumple:
    \begin{equation} \label{eq:parts}
        \int_a^b u \, dv = \Bigl[ u v \Bigr]_a^b - \int_a^b v \, du.
    \end{equation}
    \end{definition}

    \begin{lemma}[Límite fundamental] \label{lem:limite}
    Para todo entero $n \geq 0$ y todo $\lambda > 0$:
    \begin{equation}
        \lim_{x \to \infty} x^n e^{-\lambda x} = 0.
    \end{equation}
    \end{lemma}

    \begin{proof}
    Se aplica la regla de L'Hôpital $n$ veces al cociente $\dfrac{x^n}{e^{\lambda x}}$. El numerador eventualmente se convierte en una constante $n!$, mientras que el denominador siempre contiene $e^{\lambda x}$, que diverge. Por tanto, el límite es $0$.
    \end{proof}

    \section*{Demostración Paso a Paso}

    \subsection*{Paso 1: Cálculo de la esperanza $E[X]$}

    Partimos de la definición del valor esperado para una variable continua:
    $$ E[X] = \int_{0}^{\infty} x \, f(x) \, dx = \int_{0}^{\infty} x \, \lambda e^{-\lambda x} \, dx. $$

    Aplicamos integración por partes (\ref{eq:parts}) con:
    $$ u = x \quad \Rightarrow \quad du = dx, $$
    $$ dv = \lambda e^{-\lambda x} \, dx \quad \Rightarrow \quad v = -e^{-\lambda x}. $$

    Sustituyendo en (\ref{eq:parts}):
    \begin{align*}
        E[X] &= \Bigl[ -x e^{-\lambda x} \Bigr]_0^\infty - \int_{0}^{\infty} (-e^{-\lambda x}) \, dx \\[4pt]
             &= \Bigl[ -x e^{-\lambda x} \Bigr]_0^\infty + \int_{0}^{\infty} e^{-\lambda x} \, dx.
    \end{align*}

    Evaluamos el término de frontera. En el límite superior, por el Lema~\ref{lem:limite} con $n=1$, tenemos $\lim_{x \to \infty} x e^{-\lambda x} = 0$. En el límite inferior, $0 \cdot e^{0} = 0$. Por tanto, el término de frontera se anula completamente:
    $$ \Bigl[ -x e^{-\lambda x} \Bigr]_0^\infty = 0. $$

    La integral restante es directa:
    $$ \int_{0}^{\infty} e^{-\lambda x} \, dx = \left[ -\frac{1}{\lambda} e^{-\lambda x} \right]_0^\infty = 0 - \left(-\frac{1}{\lambda}\right) = \frac{1}{\lambda}. $$

    En conclusión:
    \begin{equation} \label{eq:EX}
        E[X] = \frac{1}{\lambda}.
    \end{equation}

    \subsection*{Paso 2: Cálculo del segundo momento $E[X^2]$}

    Procedemos de manera análoga:
    $$ E[X^2] = \int_{0}^{\infty} x^2 \, \lambda e^{-\lambda x} \, dx. $$

    Aplicamos integración por partes con:
    $$ u = x^2 \quad \Rightarrow \quad du = 2x \, dx, $$
    $$ dv = \lambda e^{-\lambda x} \, dx \quad \Rightarrow \quad v = -e^{-\lambda x}. $$

    Sustituyendo:
    \begin{align*}
        E[X^2] &= \Bigl[ -x^2 e^{-\lambda x} \Bigr]_0^\infty + \int_{0}^{\infty} 2x e^{-\lambda x} \, dx.
    \end{align*}

    El término de frontera se anula por el Lema~\ref{lem:limite} con $n=2$. Nos queda:
    $$ E[X^2] = \int_{0}^{\infty} 2x e^{-\lambda x} \, dx. $$

    Para conectar con la esperanza ya calculada, factorizamos y reinsertamos $\lambda$:
    \begin{align*}
        E[X^2] &= 2 \int_{0}^{\infty} x e^{-\lambda x} \, dx \\[4pt]
               &= \frac{2}{\lambda} \int_{0}^{\infty} x \, \lambda e^{-\lambda x} \, dx \\[4pt]
               &= \frac{2}{\lambda} \, E[X].
    \end{align*}

    Sustituyendo el resultado (\ref{eq:EX}):
    \begin{equation} \label{eq:EX2}
        E[X^2] = \frac{2}{\lambda} \cdot \frac{1}{\lambda} = \frac{2}{\lambda^2}.
    \end{equation}

    \subsection*{Paso 3: Cálculo de la varianza $\operatorname{Var}(X)$}

    Sustituimos los momentos (\ref{eq:EX}) y (\ref{eq:EX2}) en la fórmula (\ref{eq:varianza}):
    \begin{align*}
        \operatorname{Var}(X) &= E[X^2] - \left(E[X]\right)^2 \\[4pt]
                              &= \frac{2}{\lambda^2} - \left(\frac{1}{\lambda}\right)^2 \\[4pt]
                              &= \frac{2}{\lambda^2} - \frac{1}{\lambda^2} \\[4pt]
                              &= \frac{1}{\lambda^2}.
    \end{align*}

    \subsection*{Paso 4: Cálculo de la desviación estándar $\sigma_X$}

    Aplicando la definición (\ref{eq:sigma}):
    $$ \sigma_X = \sqrt{\operatorname{Var}(X)} = \sqrt{\frac{1}{\lambda^2}} = \frac{1}{\lambda}. $$

    Dado que $\lambda > 0$, la raíz cuadrada retorna directamente $1/\lambda$ sin necesidad de valor absoluto.

    \begin{flushright}
        \qedsymbol
    \end{flushright}

    \section*{Interpretación del Resultado}

    El cálculo nos ha llevado a que la desviación estándar de una distribución exponencial es:
    $$ \sigma_X = \frac{1}{\lambda}. $$

    Este resultado admite las siguientes interpretaciones:

    \begin{enumerate}
        \item \textbf{Coincidencia con la media.} Observamos que $\sigma_X = E[X] = \dfrac{1}{\lambda}$. Esto implica que el coeficiente de variación es exactamente $1$:
        $$ CV = \frac{\sigma_X}{E[X]} = 1. $$
        En otras palabras, la dispersión de los datos es del mismo orden de magnitud que el valor central. Esta es una propiedad distintiva de la exponencial y no ocurre en distribuciones como la normal, donde media y desviación estándar son parámetros independientes.

        \item \textbf{Un solo parámetro controla todo.} En una distribución normal $N(\mu, \sigma^2)$, existen dos parámetros: uno de localización ($\mu$) y uno de dispersión ($\sigma$). En la exponencial, el único parámetro $\lambda$ determina simultáneamente ambos aspectos. Si $\lambda$ es grande (eventos frecuentes), tanto la media como la dispersión se reducen; si $\lambda$ es pequeño (eventos raros), ambas aumentan.

        \item \textbf{Implicación práctica.} Consideremos el tiempo de espera hasta que llega un cliente a un cajero, modelado con $\lambda = 2$ clientes por minuto. La media es $E[X] = 0.5$ minutos (30 segundos) y la desviación estándar también es $\sigma_X = 0.5$ minutos. Esto significa que es perfectamente plausible observar tiempos de $1.5$ o $2$ minutos, ya que la variabilidad es tan grande como el promedio mismo. La distribución exponencial tiene una cola pesada que permite estas desviaciones significativas.

        \item \textbf{Relación con la propiedad sin memoria.} La igualdad $\sigma_X = E[X]$ es coherente con la propiedad de falta de memoria estudiada en la tarea anterior. Si la distribución no ``recordara'' cuánto tiempo ha transcurrido, no tendría sentido que la dispersión dependiera de un punto de referencia distinto al origen; el proceso se ``reinicia'' constantemente con la misma escala $1/\lambda$.
    \end{enumerate}

    \section*{Conclusiones}

    \begin{itemize}
        \item La desviación estándar de $X \sim \mathrm{Exp}(\lambda)$ es $\sigma_X = \dfrac{1}{\lambda}$.
        \item El cálculo requirió evaluar dos integrales por partes, justificando la anulación de los términos de frontera mediante límites.
        \item La varianza resultó ser $\operatorname{Var}(X) = \dfrac{1}{\lambda^2}$.
        \item La igualdad $\sigma_X = E[X]$ (coeficiente de variación unitario) revela que la dispersión es proporcional a la media, lo cual distingue a la exponencial de otras distribuciones continuas.
        \item Esta propiedad tiene consecuencias prácticas directas: modelar tiempos de espera con una exponencial implica aceptar una variabilidad inherente tan grande como el valor esperado mismo.
    \end{itemize}

    \begin{thebibliography}{}

    \bibitem{walpole2017}
    Walpole, R. E., Myers, R. H., Myers, S. L., \& Ye, K. (2017).
    \textit{Probability and Statistics for Engineers and Scientists} (9th ed.).
    Pearson.

    \end{thebibliography}

\end{document}